\section*{Chapter \ref{ch:s2:skew-and-term-structure-relative-value} --- Skew and Term-Structure Relative Value}

\begin{solution}{exo:s2:skew-and-term-structure-relative-value:1}
$100 - 10 \times (-2.3) = 123$.
\end{solution}

\begin{solution}{exo:s2:skew-and-term-structure-relative-value:2}
The vol falls by $0.4/20 = 2\%$ of itself, so the short position makes 2 units.
\end{solution}

\begin{solution}{exo:s2:skew-and-term-structure-relative-value:3}
$\ln 0.5/\ln 0.98 = 34.3$ trading days.
\end{solution}

\begin{solution}{exo:s2:skew-and-term-structure-relative-value:4}
Vol mean-reverts, so the fair term structure inverts when vol is high and slopes up when it is low: the raw slope correlates with the level at
$-0.83$ on the synthetic surface and $-0.54$ on the real VIX curve. Buying the calendar when the slope is low means buying it when vol is high, a bet
that vol falls. The residual on the level removes most of that.
\end{solution}

\begin{solution}{exo:s2:skew-and-term-structure-relative-value:5}
As a timing signal for the short-variance trades of chapter~1: sell variance when the slope says the variance premium is high. For a calendar trader
the finding says the slope is mostly the price of variance risk, not a forecast of vol, so a calendar signalled on the slope is exposed to the variance
premium, the level trade it was meant to avoid.
\end{solution}

\begin{solution}{exo:s2:skew-and-term-structure-relative-value:6}
Sell the index's put skew and buy the members' put skew (a risk reversal on the index against risk reversals on the members), vega-neutral and
delta-hedged. Its main risk is correlation: in a crash every member falls with the index and the index skew's richness is justified, which is
chapter~2's dispersion risk in skew form.
\end{solution}

\begin{solution}{exo:s2:skew-and-term-structure-relative-value:7}
The raw skew book still earns a Sharpe ratio of 1.08 (worst month $-2.8$ standard deviations): it never needed the planted effect, because it is
selling the put wing when vol is high. The residual skew book loses ($-0.41$, a worst month of $-11.7$ standard deviations), since its signal is now
only the regression's error. Both calendars lose ($-0.12$ raw, $-0.11$ residual) to costs and gamma.
\end{solution}

\begin{solution}{exo:s2:skew-and-term-structure-relative-value:8}
The attribution says otherwise. Only 9\% of that book's P\&L comes from the skew moving back, 94\% from spot and time; the raw skew is steep when vol is
high, so the book is selling the put wing when vol is high. With the planted noise switched off it still earns 1.08. The Sharpe ratio is real, but it
pays for the variance premium and the crash risk that comes with it, not for skew mean reversion.
\end{solution}

\begin{solution}{pb:s2:skew-and-term-structure-relative-value:1}
\begin{enumerate}
\item A \omterm{def:s2:skew-and-term-structure-relative-value:skew}{skew trade}: long and short options at two strikes of one expiry, vega-neutral and delta-hedged, on the difference in their vols. A
\omterm{def:s2:skew-and-term-structure-relative-value:term}{term-structure trade}: the same across two maturities.
\item $100 - 10S$, with $S$ the price of S\&P 500 skewness; 123.2 on average since 1990, 110.0 to 147.3 for the middle 90\%, 101.23 at the low
(19 March 1991) and 183.12 at the high (18 February 2025).
\item Individual stocks' risk-neutral distributions are far less negatively skewed than the index's.
\item The fair skew is $-0.12$ at the average level and moves by $-0.4$ per unit of one-month vol; the noises are an AR(1) on the skew (sd 0.06) and on
the three-month vol (1.5 vol points), with a half-life of 34 days.
\item 1.93 points on average; inverted on 7.6\% of days; $-18.23$ on 12 March 2020.
\item Vol mean-reverts, so high vol inverts the curve; the correlation is $-0.54$.
\item The slope carries the price of variance risk and predicts variance-swap, VIX-future and straddle returns.
\item RVX 5.2 points above VIX (sd 2.3, half-life 15.7 days); VXN 3.0 above (sd 2.3, half-life 15.3).
\item Level, slope and skew from the surface; slope and skew regressed on the level with expanding windows, keeping the residuals.
\item One unit per one-percent rise in the leg's vol at inception; P\&L in percent of vol.
\item Repricing in order: spot and time, then level, then skew noise, then slope noise; the buckets add up, and moving the level first puts the fair
skew and slope changes in the level bucket.
\item 1.13, 0.75, 0.07 and 0.64; spot and time 94, 73, $-31$ and 97\%; level 5, 24, 160 and $-16$\%; skew or slope 9, 16, 155 and 51\%.
\item \textbf{Named result.} The skew and calendar trades signalled on residuals earn Sharpe ratios of 0.75 and 0.64 on this seed (0.39 and 0.38 on
average over eight). Only 16\% and 51\% of their P\&L comes from the planted skew and slope; the raw skew book earns 1.13 with 9\%, and the raw
calendar's P\&L is 160\% level.
\item Sharpe ratios of 1.01 (skew bucket) and 1.51 (slope bucket).
\item Means of 0.90, 0.39, $-0.07$ and 0.38; this seed is on the lucky side for the residual books.
\item The risk reversal: 12.0 units a month from spot and time, $-15.9$ from the level, $-6.2$ in all. The calendar: 9.6, $-3.2$, 4.7.
\item Long two units of call wing into a 4.5\% rally with the one-month vol falling from 15.2\% to 9.7\%: $-289$ units, 4.5 standard deviations.
\item The raw skew book still earns 1.08; the residual books and calendars lose.
\item Attribute the P\&L, compare raw and residual signals, include joint moves of spot and vol, charge option costs, report several seeds.
\item Whatever its legs do not match: the joint moves of spot and vol, and near-month gamma.
\end{enumerate}
\end{solution}

\begin{solution}{iq:s2:skew-and-term-structure-relative-value:1}
Long an out-of-the-money call and short an out-of-the-money put (or the reverse) at equal deltas or vegas. Its price, quoted as the call's vol minus
the put's, measures the skew: a large negative value says the put wing is dear.
\end{solution}

\begin{solution}{iq:s2:skew-and-term-structure-relative-value:2}
Vol rises and the curve inverts: the short one-month leg's vol rises more than the long three-month's, so the level move loses; the short near leg has
more gamma than the long far leg, so the large moves lose too. Both hurt.
\end{solution}

\begin{solution}{iq:s2:skew-and-term-structure-relative-value:3}
Regress the skew on the level and check the correlation; trade the residual as well as the raw signal and compare; attribute the P\&L into level and
skew buckets. If the raw signal's P\&L is mainly spot, time or level, it is a level signal.
\end{solution}

\begin{solution}{iq:s2:skew-and-term-structure-relative-value:4}
Gamma and theta that do not cancel; vega that appears as spot moves (a short wing coming to the money); the joint move of spot, vol and skew; jumps
between hedges; basis between underlyings; the costs of rolling and hedging.
\end{solution}

\begin{solution}{iq:s2:skew-and-term-structure-relative-value:5}
Reprice each position in a fixed order of market moves (spot and time, then level, then shape parameters) and record each step's change; the last
state must equal the next day's full surface, so the buckets add to the total by construction. The surface must be a function of a small set of
parameters, with a rule for how it moves with spot.
\end{solution}

\begin{solution}{iq:s2:skew-and-term-structure-relative-value:6}
Over $\mathrm{d}t$ the hedged option's value changes by $\Theta\,\mathrm{d}t + \tfrac12\Gamma(\mathrm{d}S)^2$, and Black--Scholes gives $\Theta =
-\tfrac12\Gamma S^2\sigma_i^2$; with $(\mathrm{d}S)^2 \approx S^2\sigma_r^2\,\mathrm{d}t$ the P\&L is $\tfrac12\Gamma S^2(\sigma_r^2 -
\sigma_i^2)\,\mathrm{d}t$. Vega is $S\varphi(d_1)\sqrt\tau$ and gamma $\varphi(d_1)/(S\sigma_i\sqrt\tau)$, so gamma per unit of vega is
$1/(S^2\sigma_i\tau)$: the leg with the higher implied vol has less gamma per unit of vega.
\end{solution}
